\documentclass[11pt]{article}
\usepackage[margin=1in]{geometry}
\usepackage{amsmath,amssymb}
\usepackage{booktabs}
\usepackage{mathrsfs}

\title{Surface Roller Model for \texttt{WAVE\_OFFLINE} Coupling\\[4pt]
\large Porting from WKB to Offline WEC}
\author{}
\date{}

\begin{document}
\maketitle

\section{What is the surface roller?}

When a wave breaks, not all of the breaking energy is instantly transferred to
the mean current. A fraction $\alpha_r$ of the breaking dissipation
$\varepsilon^b$ feeds a \emph{surface roller}---a mass of aerated water that
rides on the front face of the broken wave and propagates shoreward at the
wave phase speed~$c$ (not the group speed~$c_g$). As the roller propagates,
it dissipates at a rate~$\varepsilon^r$ that depends on the roller energy and
wave properties. It is $\varepsilon^r$, not~$\varepsilon^b$, that ultimately
accelerates the mean nearshore current and drives breaking-induced turbulent
mixing.

The physical consequence is a \textbf{shoreward displacement} of the effective
momentum source: breaking occurs at the bar crest, but the roller carries that
energy shoreward before releasing it. This broadens the cross-shore profile of
the forcing and shifts the peak $v$ location toward the shore (U2010, Fig.~6).

\section{Roller energy balance}

The roller action density is defined analogously to the primary wave action:
\begin{equation}
  \mathscr{A}^r = \frac{E^r}{\sigma},
  \label{eq:roller_action}
\end{equation}
where $E^r$ is the roller energy density and $\sigma$ is the intrinsic wave
frequency. Its evolution equation (U2010 Eq.~40) is
\begin{equation}
  \frac{\partial\mathscr{A}^r}{\partial t}
  + \nabla\!\cdot\!\bigl(\mathscr{A}^r\,\mathbf{c}\bigr)
  = \frac{\alpha_r\,\varepsilon^b - \varepsilon^r}{\sigma},
  \label{eq:roller_balance}
\end{equation}
where $\mathbf{c} = (\sigma/k^2)\,\mathbf{k}$ is the phase velocity vector
(not the group velocity), and the two source/sink terms are:

\begin{itemize}
\item \textbf{Source:} $\alpha_r\,\varepsilon^b/\sigma$. A fraction
  $\alpha_r \in [0,1]$ of the primary wave breaking dissipation is converted
  into roller energy. The remainder $(1-\alpha_r)\varepsilon^b$ acts
  directly as breaking acceleration. Following U2010 Table~1,
  $\alpha_r = 1.0$ (all breaking feeds the roller) is the baseline; smaller
  values allow part of $\varepsilon^b$ to bypass the roller.

\item \textbf{Sink:} $\varepsilon^r/\sigma$. The roller dissipates at a rate
  (U2010 Eq.~41)
  \begin{equation}
    \varepsilon^r = \frac{g\sin\beta\;E^r}{c}
                  = g\sin\beta\;\mathscr{A}^r\,k,
    \label{eq:roller_dissip}
  \end{equation}
  where $\sin\beta \approx 0.1$ is an empirical roller slope parameter and
  $c = \sigma/k$ is the phase speed. The dissipation rate is proportional to
  the roller energy and to the local wavenumber---short, steep waves dissipate
  rollers faster than long, gentle ones.
\end{itemize}

\subsection*{Steady-state estimate}

In a steady state ($\partial\mathscr{A}^r/\partial t = 0$) with weak
advective divergence, the local balance gives
\begin{equation}
  \varepsilon^r \approx \alpha_r\,\varepsilon^b,
\end{equation}
i.e.\ the roller dissipation equals the source. The \emph{total} effective
dissipation entering the momentum equations is
$\varepsilon^{\mathrm{tot}} = (1-\alpha_r)\varepsilon^b + \varepsilon^r
\approx \varepsilon^b$, confirming that the roller redistributes, but does
not create or destroy, momentum. Its value lies in the spatial shift: the
$\varepsilon^r$ peak is displaced shoreward of the $\varepsilon^b$ peak by
the advection term $\nabla\!\cdot(\mathscr{A}^r \mathbf{c})$.

\section{Modified breaking stress (U2010 Eq.~48)}

Without the roller, the depth-averaged breaking acceleration is
\begin{equation}
  \bar{\mathbf{B}}^b = \frac{\varepsilon^b}{\rho_0\,\sigma\,D}\,\mathbf{k}.
  \label{eq:brk_noroller}
\end{equation}
With the roller, the momentum source becomes the \emph{combined}
dissipation (U2010 Eq.~50):
\begin{equation}
  \boxed{
  \bar{\mathbf{B}}^b
  = \frac{(1-\alpha_r)\,\varepsilon^b + \varepsilon^r}
         {\rho_0\,\sigma\,D}\;\mathbf{k}.
  }
  \label{eq:brk_roller}
\end{equation}
In the code, this means replacing \texttt{ebrk} with
$\texttt{(1-alpha\_r)*ebrk + rdsp}$ in the computation of
\texttt{brk(i,j)}. With $\alpha_r=1$, the direct breaking term vanishes
entirely and all momentum comes through $\varepsilon^r$.

The 3D body force version (U2010 Eq.~48) includes a vertical shape
function~$f^b(z)$:
\begin{equation}
  \mathbf{B}^b(z) = \frac{(1-\alpha_r)\,\varepsilon^b + \varepsilon^r}
                         {\rho_0\,\sigma}\;\mathbf{k}\,f^b(z),
  \qquad \int_{-h}^{\zeta} f^b(z')\,dz' = 1.
  \label{eq:brk3d_roller}
\end{equation}

\section{Modified breaking-induced mixing (U2010 Eq.~59)}

The breaking-induced eddy viscosity $K_v^b$ (added to KPP's $K_v$) is
\begin{equation}
  \boxed{
  K_v^b = c_b\left(\frac{(1-\alpha_r)\,\varepsilon^b + \varepsilon^r}
                       {\rho_0}\right)^{\!1/3}\! H_*\,D\;f^{K_v}(z),
  }
  \label{eq:Akb_roller}
\end{equation}
where $c_b \approx 0.03$ (U2010 Table~2, calibrated for DUCK94),
$H_* = a_{K_v}\,H_s$ is a length scale based on the wave height, and
$f^{K_v}(z)$ is a shape function (Type~II in the code, decaying from
surface). The exponent $1/3$ comes from the TKE-based scaling
$K \sim \ell\,q \sim H_*\,({\varepsilon}/{\rho_0})^{1/3}$.

In the code, this means replacing \texttt{ebrk} with
$\texttt{(1-alpha\_r)*ebrk + rdsp}$ inside the Akb computation block.
The shape function and $c_b$ parameter are unchanged.

\section{Roller Stokes drift}

The roller also contributes to the Stokes drift. The roller Stokes
transport (U2010 Eq.~42) is
\begin{equation}
  \mathbf{U}^r = \frac{E^r}{\rho_0\,\sigma}\,\mathbf{k}
               = \frac{\mathscr{A}^r}{\rho_0}\,\mathbf{k},
  \label{eq:roller_stokes}
\end{equation}
so the total Stokes transport (U2010 Eq.~43) is
\begin{equation}
  \mathbf{U}^{\mathrm{St}} = \frac{E + E^r}{\rho_0\,\sigma}\,\mathbf{k}
  = \frac{\mathscr{A} + \mathscr{A}^r}{\rho_0}\,\mathbf{k}.
\end{equation}
The vertical distribution is assumed to follow the same profile as the
primary Stokes drift (U2010 Eq.~44). In the code, this means multiplying the
existing Stokes drift by $(1 + \mathscr{A}^r/\mathscr{A})$ or, equivalently,
replacing the wave action $\mathscr{A}$ with
$\mathscr{A} + \mathscr{A}^r$ in the Stokes drift computation.

\section{Advection scheme}

The roller propagates at the wave \textbf{phase speed}~$c = \sigma/k$,
not the group speed. In the WKB code (\texttt{wkb\_wwave.F}), the advection
uses a 1st-order upwind scheme:
\begin{align}
  c^u_{i+1/2} &= \tfrac{1}{2}\bigl(c^x_{i+1,j} + c^x_{i,j}\bigr), \\
  c^v_{j+1/2} &= \tfrac{1}{2}\bigl(c^\eta_{i,j+1} + c^\eta_{i,j}\bigr),
\end{align}
where $c^x = c\cos\theta$, $c^\eta = c\sin\theta$ are phase velocity
components in the wave propagation direction. The upwind fluxes are
\begin{align}
  \mathcal{F}^x_{i+1/2} &= \tfrac{1}{2}\,\Delta n_u\bigl[
    (c^u - |c^u|)\,\mathscr{A}^r_{i+1} + (c^u + |c^u|)\,\mathscr{A}^r_{i}
  \bigr], \\
  \mathcal{F}^\eta_{j+1/2} &= \tfrac{1}{2}\,\Delta m_v\bigl[
    (c^v - |c^v|)\,\mathscr{A}^r_{j+1} + (c^v + |c^v|)\,\mathscr{A}^r_{j}
  \bigr].
\end{align}
The update is
\begin{equation}
  \mathscr{A}^r_{i,j}^{n+1} = \mathscr{A}^r_{i,j}^{n}
  - \Delta t\,p_m\,p_n\bigl(
    \mathcal{F}^x_{i+1/2} - \mathcal{F}^x_{i-1/2}
    + \mathcal{F}^\eta_{j+1/2} - \mathcal{F}^\eta_{j-1/2}
  \bigr)
  + \Delta t\bigl(\alpha_r\,\varepsilon^b/\sigma
                 - g\sin\beta\;\mathscr{A}^r\,k/\sigma\bigr).
  \label{eq:upwind_update}
\end{equation}

\subsection*{CFL constraint}

The upwind scheme is stable for $\mathrm{CFL} = c\,\Delta t/\Delta x \leq 1$.
In shallow water, $c \to \sqrt{gD}$. Representative values for
$\Delta x = 30$\,m, $\Delta t = 3$\,s:

\medskip
\begin{center}
\begin{tabular}{@{}cccc@{}}
\toprule
$D$ (m) & $c$ (m/s) & CFL & Stable? \\
\midrule
2  & 4.4  & 0.44 & yes \\
5  & 7.0  & 0.70 & yes \\
10 & 9.9  & 0.99 & marginal \\
20 & 14.0 & 1.40 & \textbf{no} \\
\bottomrule
\end{tabular}
\end{center}

\medskip\noindent
For L4PTSAL, the surfzone (where rollers exist) typically has $D < 10$\,m,
so CFL $< 1$ is satisfied. In deeper water, $\varepsilon^b \to 0$, so
the source term vanishes and $\mathscr{A}^r$ decays exponentially---the
advection there is irrelevant. If needed, a CFL limiter or subcycling can
be added for safety.

\section{Connection to WW3 forcing}

With \texttt{WAVE\_OFFLINE}, the wave fields ($H_s$, $\sigma$, $\theta$,
$\varepsilon^b$, $\varepsilon^{wd}$, $Q_b$) are read from a WW3 forcing
file and interpolated to each baroclinic time step. For the roller:

\begin{itemize}
\item $\varepsilon^b$ (\texttt{wveb}): already read. This is the source
  term for the roller.
\item $Q_b$ (\texttt{wvqb}): fraction of breaking waves. Already declared
  in the code under \texttt{SURFACE\_ROLLER} with the forcing reader
  (\texttt{nc\_wwqb}). Available in the forcing file.
  Used for diagnostic initialization (see below), not needed for the
  prognostic equation itself.
\item $k$, $\sigma$ (\texttt{kw}, \texttt{fr}): computed from dispersion
  relation. Phase speed $c = \sigma/k = \texttt{fr/kw}$.
\item $\theta$ (\texttt{wdrx}, \texttt{wdre}): wave direction components.
  Phase velocity vector: $c^x = c\,\texttt{wdrx}$,
  $c^\eta = c\,\texttt{wdre}$.
\end{itemize}

\subsection*{Diagnostic initialization using $Q_b$}

The roller can be initialized from a diagnostic steady-state estimate
(Svendsen, 1984; U2010 Eq.~45):
\begin{equation}
  E^r_{\mathrm{diag}} = \rho_0\,g\,D\,A_R\,Q_b,
\end{equation}
where $A_R = a_R^0\,H_s\,L_p$ is the roller area with empirical constants
$a_R^0 \approx 0.06$ and $L_p = 2\pi/|\mathbf{k}_p|$ is the peak
wavelength. This avoids the cold-start transient from $\mathscr{A}^r = 0$.

\section{Relationship to \texttt{SMOOTH\_BRK}}

\texttt{SMOOTH\_BRK} (Laplacian smoothing of \texttt{brk}) is a
\emph{surrogate} for the roller: it spreads the breaking momentum
isotropically over a Gaussian width $\sigma = \Delta x\sqrt{3N/8}$.
The roller is the \emph{physical} mechanism:

\medskip
\begin{center}
\begin{tabular}{@{}lcc@{}}
\toprule
Property & \texttt{SMOOTH\_BRK} & Roller \\
\midrule
Direction        & isotropic & shoreward (along $\mathbf{k}$) \\
Spreading width  & tunable ($N$) & physical ($\beta$, $c$) \\
Conservation     & exact (Neumann BC) & exact (PDE) \\
Akb modification & no & yes (Eq.~\ref{eq:Akb_roller}) \\
Stokes drift     & no & yes (Eq.~\ref{eq:roller_stokes}) \\
CFL constraint   & none & $c\,\Delta t/\Delta x < 1$ \\
Complexity       & $\sim$100 lines & $\sim$300 lines \\
\bottomrule
\end{tabular}
\end{center}

\medskip\noindent
With the roller implemented, \texttt{SMOOTH\_BRK} becomes redundant and
should be disabled (\texttt{\#undef SMOOTH\_BRK}). The two should not be
used simultaneously, as that would double-smooth the forcing.

\end{document}
